%%%%

%%%   Самарский государственный технический университет

%%%   Кафедра прикладной математики и информатики

%%%

%%%   Преамбула для статьи в журнал "Вестник Самарского государственного

%%%   технического университета. Серия: «Физико-математические науки»"

%%%   Ориентирована под MikTEx 2.4 for Win32

%%%

%%%   Хотя сам журнал собирается под GNU/Linux на дистрибутиве TeXLive



\documentclass[11pt,a4paper,twoside]{article}



\usepackage[cp1251]{inputenc}

\usepackage[T2A]{fontenc}

\usepackage[english,russian]{babel}

\usepackage{samgtubib}

\usepackage[dvips]{graphicx}

\usepackage[multiple]{footmisc}



\usepackage{samgtu}

\renewcommand{\VestnikYear}{2009} % Год издания Вестника

\renewcommand{\VestnikNumber}{2\,(19)} % Номер Вестника

\renewcommand{\VestnikMonth}{Ноябрь} % Месяц выхода Вестника

\renewcommand{\VestnikData}{01.11.09} % Дата подписания Вестника в печать

\renewcommand{\VestnikUPL}{26,64} % Усл. печ. л.

\renewcommand{\VestnikUIL}{25,73} % Уч.-изд. л.

\renewcommand{\VestnikRegNumber}{479/09} % Регистрационный номер

\renewcommand{\VestnikOrder}{994} % Номер заказа







%

%

%\usepackage{nccboxes}

%\begin{document}



%\newcommand{\totop}[1]{\raisebox{6pt}[1pt][2pt]{#1}}

%\newcommand{\tobot}[1]{\raisebox{-6pt}[1pt][2pt]{#1}}



\renewcommand{\firstpage}{173}

\renewcommand{\lastpage}{177}





\udk{534.222.2} \msc{76N15}





\title{Расч\"eт метательной способности взрывчатых веществ при цилиндрическом и~торцевом метаниях металла}

{Расч\"eт метательной способности взрывчатых веществ \dots}



\author{И.\,И.~Реут, А.\,Л.~Кривченко}{Р\,е\,у\,т~И.\,И., К\,р\,и\,в\,ч\,е\,н\,к\,о~А.\,Л.}



\institute{\samgtu.}



\email{E-mails}{enterfax@mail.ru, snex@rambler.ru}



\otherinf{\fullauthor{Реут Игорь Игоревич}  \position{аспирант} \dept{каф. защиты в чрезвычайных ситуациях}

\fullauthor{Кривченко Александр Львович} \regalia{д.т.н., доцент} \position{профессор} \dept{каф. защиты в чрезвычайных ситуациях}}



\keywords{метательная способность, торцевое метание, цилиндрическое расширение оболочки}



\workabstract{Предложены методики расчета относительной скорости торцевого метания медной пластины методом эквивалентных

масс и скорости расширения медной цилиндрической оболочки.

Результаты вычислений сопоставлены с экспериментальными данными.

Среднеквадратическая ошибка составила менее 3,5\,{\%}.}



\engtitle{Computing of throwing ability of the explosives at cylindric front tossing of metal}



\engauthor{I.\,I.~Reout, A.\,L.~Krivchenko}{R\,e\,o\,u\,t~I.\,I., K\,r\,i\,v\,c\,h\,e\,n\,k\,o~A.\,L.}





\enginstitute{\engsamgtu.}



\otherenginfm{\fullauthor{Igor I. Reout}

\position{Postgraduate Student} \dept{Dept. of Protection of Extremal Situation} \fullauthor{Alexander~L.~Krivchenko} \regalia{Dr. Sci. (Techn.)} \position{Professor} \dept{Dept. of Protection of Extremal Situation}}





\engabstract{Calculation methods of copper plate relative front tossing speed by equivalent mass method and the widening speed calculation methods for the copper cylindric confinement are offered. The calculation results were compared with experimental data. The mean-square deviation is less than 3,5 percent.

}



\engkeywords{throwing ability, front tossing, cylindric shell extension}



\date{13/V/2011}

\revisiondate{13/VI/2011}





\maketitle



\small



\Section[n]{Введение} Сложные динамические взаимодействия

представляют самостоятельный интерес для различных областей науки

и техники, в том числе и для разработки новых систем динамического

оружия. Существенное место среди кумулятивных зарядов занимают так

называемые снарядоформирующие заряды, формирующие

высокоскоростные компактные или удлинённые поражающие элементы,

называемые также ударными ядрами.



В настоящее время не существует теоретических методов оценки эффективности

действия сплавов для облицовки боеприпасов типа <<ударное ядро>>. Подбор

материалов ограничен экспериментальными данными в основном по меди и стали.



В отличие от таких параметров, как скорость детонации и давление продуктов

взрыва за фронтом, величина метательной способности не имеет абсолютного

значения, а является относительной характеристикой в выбранном виде

испытания. Удобно метательную способность выражать в относительных единицах.

Относительная метательная способность, как правило, определятся по

некоторому эталону взрывчатого вещества (ВВ) и выражается в процентах по

отношению к соответствующей характеристике эталона.



\smallskip

\Section[n]{Расчёт относительной скорости

торцевого метания на основе модели детонационной оптики} 

В~общем виде \cite{reut1:bib1} скорость разгона пластины (метаемого элемента) $W= 2U$, где $U$ "--- массовая скорость продуктов детонации. 

Это выражение, полученное методом отражения, можно заменить соотношением для эквивалентных масс метаемой пластины и заряда.

При этом оценка скорости метаемой пластины по методу торцевого метания требует знания активной массы~ВВ~\cite{reut1:bib2}.



Нами предложен метод эквивалентных масс, в котором энергетика ВВ, участвующего в~разгоне пластины, сопоставляется с~метаемым элементом в~соотношении~1:1. В~качестве критерия принимается эквивалентное взаимодействие ВВ с~метаемым элементом.



Известно~\cite{reut1:bib9}, что при постоянной  высоте ВВ его энергия $E$, участвующая в метании, определяется соотношением

$$

E=m_0U^2,

$$

где $

 m_0= \pi r_0^2 h_0 \rho_0/3,

$ "--- активная масса ВВ, участвующая в разгоне пластины и занимающая объём конуса с радиусом основания $r_0$  и высотой $h_0$ \cite{reut1:bib2};

$\rho_0$ "--- плотность заряда ВВ. 

При этом метаемая круглая пластина радиуса $r_{\text{пл}}$  с толщиной $h_{\text{пл}}$, имеющая массу $m_{\text{пл}}= \pi r_{\text{пл}} h_{\text{пл}}\rho_{\text{пл}}$ ($\rho_{\text{пл}}$ "--- плотность метаемой пластины), получает энергию

$$

 E_\text{пл} =\frac{m_\text{пл} W^2}{2}.

$$

Приравнивая $E$ и $E_{\text{пл}}$ и учитывая, что $r_0=r_{\text{пл}}$, получим

\begin{equation*}

\frac{W^2}{U^2}= \frac 23 \frac{h_0\rho _0 }{h_\text{пл}\rho _\text{пл}},

\end{equation*}

т.\,е. скорость метания при эквивалентных массах заряда и пластины определяется из выражения

\begin{equation}

\label{reut1:eq6}

W_\text{экв} =U\sqrt {\frac 23 \frac{h_0\rho _0 }{h_\text{пл}\rho _\text{пл}}}.

\end{equation}





\begin{table}[b!] \centering \small

\setlength{\tabcolsep}{3.5pt}

\caption{\bf Данные относительной скорости метания

по~методике М--60} \vspace{-3mm}

\begin{tabular}{l|c|c|c|c} \hline

\multicolumn{1}{c|}{\footnotesize  ВВ} & 

\footnotesize $\rho_0$, г/см$^{3}$ & 

\footnotesize $W_{\text{эксп}}$, {\%} & 

\footnotesize $W_\text{экв}$, {\%} & 

\footnotesize {\%} ош. \\ \hline

\rule{0mm}{12pt}%

Октоген   & 1,87 & 100,0 & 100,0 & --- \\ 

Гексоген  & 1,76 & \phantom{1}97,0& \phantom{1}94,4& $-2,7$ \\ 

ТНТ       & 1,57 & \phantom{1}76,0& \phantom{1}73,2& $-3,7$ \\ 

Бис-(фтординитроэтил)-нитрамин & 1,87 & \phantom{1}93,7 & \phantom{1}98,0 & $+4,6$ \\ 

Бис-(хлординитроэтил)-нитрамин & 1,88& \phantom{1}90,3 & \phantom{1}95,1 & $+5,3$ \\ 

1,4-динитропиперазин & 1,58 & \phantom{1}77,2 & \phantom{1}71,4 & $-4,5$ \\ 

2',2',2'-тринитроэтил-4,4,4-тринитробутират & 1,75 & \phantom{1}93,3 & \phantom{1}97,1 & $+4,1$ \\ 

Бис-(тринитроэтил)-формаль & 1,65 & \phantom{1}87,7 & \phantom{1}91,1 & $+3,9$ \\ 

ТЭН & 1,74 & \phantom{1}93,7 & \phantom{1}97,7 & $+4,2$ \\ 

ГНБ & 1,97 & 109,7& 115,8& $+5,6$ \\ 

ТНБ & 1,64 & \phantom{1}79,4 & \phantom{1}80,0 & $+0,7$ \\ 

ТАТБ& 1,85 & \phantom{1}82,5 & \phantom{1}78,8 & $-4,5$ \\ 

Тетрил & 1,71 & \phantom{1}85,7 & \phantom{1}86,1 & $+0,5$ \\ 

Гексанитростильбен & 1,64 & \phantom{1}79,0 & \phantom{1}79,1 & $+0,1$ \\ 

3,3'-динитро-4,4'-азоксифуразан & 1,63 & \phantom{1}77,0 & \phantom{1}76,6 & $-0,5$ \\ 

3,3'-динитро-4,4'-азоксифуразан & 1,79 & 101,7& 107,1& $+5,3$ \\ 

Динитро-фуроксанил & 1,92 & 109,0 & 111,2 & $+2,0$ \\ 

3,3'-диамино-4,4'-азоксифуразан & 1,72 & \phantom{1}96,9 & \phantom{1}96,0 & $-1,0$ \\ 

Бензотрифуроксан & 1,88 & 101,5 & 104,3& $+2,7$ \\ 

ДНТФ &1,90&103,8&108,8&+4,8 \\ 

7-амино-4,6-динитробензофуроксан & 1,79 & \phantom{1}88,2 & \phantom{1}87,6 & $-0,7$ \\ \hline

\end{tabular}

\end{table}



В случае неэквивалентного метания, когда масса пластины больше или меньше массы заряда, скорость метания соответственно изменится. 

Данный метод может быть предложен как наиболее удобный и~обобщённый.



По разработанному методу эквивалентных масс можно рассчитать относительную метательную способность некоторых ВВ и~сопоставить

с~экспериментальными данными торцевого метания. 

В~табл.~1 представлены экспериментальные и расчётные данные относительной скорости метания медной пластины ($\rho_{\text{пл}}=8,93$~г/см$^3$) по методике М--60 [\citen{reut1:bib3}--\citen{reut1:bib5}]. Здесь расчётные данные скорости метания пластины были

получены с~использованием выражения \eqref{reut1:eq6} при $h_0=h_{\text{пл}}$.







\smallskip



\Section[n]{Расчёт скорости расширения медной цилиндрической оболочки}

При разработке теоретического метода оценки метательной способности

индивидуальных ВВ по методике Т--20 использовался тот же подход, что и в \cite{reut1:bib6},

т.\,е. предполагалось, что кинетическая энергия ускоряемого тела пропорциональна объёмному энергосодержанию ВВ, а~коэффициент отбора энергии зависит от объёмного числа молей газообразных продуктов взрыва.



Известно, что распределение скорости звука и массовой скорости частиц в~плоской детонационной волне линейно. 

При этом состояние точки Жуге определяется по формуле $D=C+U$, где $D$ "--- скорость ударной волны;

$C$ "--- объёмная скорость звука.

%Для разработанного метода оценки скорости расширения оболочки был использован массив, состоящий из 25~ВВ. 

Нами предлагается уравнение по определению скорости расширения медной цилиндрической оболочки в виде

$$

 W_{\text{цил}} =({C_{\text{ВВ}} +U})/{2,66},

$$

где делитель численно связан с коэффициентами уравнения ударной адиабаты $D=\alpha + \lambda U$ соотношением $\alpha/\lambda$, т.\,е. $3,98/1,495=2,66$; $C_{\text{ВВ}}$ "--- объёмная скорость звука~в~ВВ.





\begin{table}[b!]   \vspace{-3mm}

\setlength{\tabcolsep}{4pt}

\centering \small

\caption{\bf Данные скорости расширения медной цилиндрической оболочки по методике Т--20} \vspace{-3mm}

\begin{tabular}{l|c|c|c|c} \hline

\multicolumn{1}{c|}{\tobot{\footnotesize ВВ}} & 

\tobot{\footnotesize $\rho_0$, г/см$^{3}$} & 

\multicolumn{2}{c|}{\footnotesize $ W_{\text{цил}}$, м/с} & 

\tobot{\footnotesize {\%} ош.} \\ \cline{3-4}

    &                      &\footnotesize ~~~~~~~расчёт~~~~~~~& \footnotesize эксперимент & \\ 

\hline

\rule{0mm}{12pt}%

Нитрометан & 1,14&1101&1092&$+0,8$ \\ 

Бис-(нитраминометил)-нитрамин &1,85&1965&1907&$+3,0$ \\

Гексоген& 1,80&1779&1781&$-0,1$ \\ 

Бис-(тринитроэтил)-нитрамин&1,85&1541&1530&$+0,7$ \\ 

Бис-(тринитроэтил)-нитрамин&1,90&1607&1650&$-2,6$ \\ 

Бис-(тринитроэтил)-нитрамин&1,94&1781&1792&$-0,6$ \\ 

Октоген&1,73&1654&1660&$-0,4$ \\ 

Октоген&1,81&1796&1770&$+1,4$ \\ 

Октоген&1,89&1853&1840&$+0,7$ \\ 

Октоген&1,91&1989&1910&$+4,1$ \\ 

Октанитро-диазаоктан&1,86&1771&1683&$+5,2$ \\ 

Тринитроэтиловый эфир &\tobot{1,78}&\tobot{1548}&\tobot{1507}&\tobot{$+2,7$} \\ 

тринитромасляной кислоты & & & & \\ 

Метилнитрат&1,21&1170&1133&$+3,2$ \\ 

Диэтанолнитраминдинитрат&1,67&1678&1616&$+3,8$ \\ 

ТЭН&1,67&1552&1661&$-6,6$ \\ 

ТЭН&1,76&1704&1790&$-4,8$ \\ 

ГНБ&1,96&1726&1710&$+0,9$ \\ 

ГНБ&1,86&1582&1616&$-2,1$ \\ 

ТАТБ&1,80&1542&1470&$+4,9$ \\ 

ТНТ&1,63&1470&1460&$+0,7$ \\ 

Тетрил&1,70&1527&1594&$-4,2$ \\ 

Гексанитростильбен&1,65&1389&1440&$-3,6$ \\ 

5-оксо-3-нитро-1,2,4-триазол&1,85&1302&1330&$-2,1$ \\ 

Окфол-3,5&1,78&1710&1750&$-2,3$ \\ 

A-IX-1&1,67&1697&1650&$+2,8$ \\ \hline

\end{tabular}

\end{table}





Для оценки предложенного расчётного метода определения $ W_{\text{цил}}$ целесообразно сопоставить

результаты вычислений с экспериментальными данными. 

Расчётные и экспериментальные  значения скорости метания приведены в~табл.~2.

Экспериментальные данные брались из работ [\citen{reut1:bib3}--\citen{reut1:bib5}, \citen{reut1:bib7}].

Расчёт $C_{\text{ВВ}}$ проводился по методу Рао \cite{reut1:bib8}, массовая скорость $U$

вычислялась по формуле \cite{reut1:bib9}

$$

U=\sqrt {\frac{\alpha_{\text{к}}+1,65}{5,5}\rho _0 Q_{\text{м}}},

$$

где $\alpha_{\text{к}}$ "--- кислородный коэффициент ВВ; $Q_{\text{м}}$ "--- максимальная теплота (энергия)

взрыва, которая отвечает протеканию наиболее теплотворных реакций при разложении ВВ и~соответствует максимуму энтропии.







\smallskip



\Section[n]{Анализ результатов и выводы}

Как видно из табл.~1 и~2, относительно высокие погрешности в~расчёте скорости метания, по"=видимому, говорят о том, что устойчивая стационарная детонация, самопроизвольно распространяющаяся со скоростью, постоянной для данного вещества, происходит при условии, когда скорость детонационной волны относительно продуктов реакции выше скорости звука в них. 

Это обстоятельство указывает на то, что, возможно, во время эксперимента с~помощью мощной ударной волны  в среде была возбуждена  детонация с большей скоростью.

Так, в работе \cite{reut1:bib10} предложен один из возможных способов увеличения коэффициента отбора энергии от ВВ за счёт пересжатой детонационной волны, распространяющейся по заряду. 

Как показали модельные расчёты, в~этом случае можно ожидать увеличения коэффициента отбора энергии от ВВ на $10$\,{\%} 

(в~зависимости от соотношения массы металла и~ВВ). 

Суть предполагаемого эффекта заключается в~том, что за счёт пересжатия детонационной волны продукты взрыва за её фронтом имеют большую удельную кинетическую энергию, а потому передают оболочке б{\'о}льший импульс, чем в~случае выхода на стационарный режим детонации. 

Таким образом, пересжатая волна, возникающая в~заряде ВВ, характеризуется увеличением сжатия и~массовой скорости во фронте детонационной волны по мере её распространения, при этом происходит увеличение коэффициента преобразования химической

энергии ВВ в~энергию потока продуктов взрыва.



Следует отметить, что разработанные методики расчёта дают малую среднеквадратичную ошибку.

Так, среднеквадратичная ошибка определения $W_\text{экв}$ по методу эквивалентных масс составляет 3,33\,{\%}, а среднеквадратичная  ошибка определения $ W_{\text{цил}}$ составляет 1,73\,{\%}.



% \smallskip

%

%\Section{Выводы}

%Предложены методики расчёта:

%

%"--- относительной скорости торцевого метания методом эквивалентных масс,

%разработанным на основе модели детонационной оптики;

%

%"--- скорости расширения медной цилиндрической облицовки.

%

%Разработанные теоретические методики прогнозирования скоростей метания по

%методам Т--20 и М--60 дают среднеквадратическую ошибку менее 3,5{\%}.



\begin{thebibliography}{99}



\RBibitem{reut1:bib1}

\by Альтшулер~Л.\,В.

\paper Применение ударных волн в физике высоких давлений

\jour УФН

\yr 1965

\vol 85

\issue 2

\pages 197--258

\transl англ. пер.:~

\by Al'tshuler~L.\,V.

\paper Use of Shock Waves in High-Pressure Physics

\yr 1965 

\jour Sov. Phys. Usp.

\vol 8 

\issue 1

\pages 52--91

% doi:10.1070/PU1965v008n01ABEH003062



\RBibitem{reut1:bib2}

\by Андреев~С.\,Г.,  Бабкин~А.\,В., Баум~Ф.\,А. и др.

\book Физика взрыва

\bookinfo в 2-x томах

\ed Л.~П.~Орленко

\bookvol 1

\publaddr М.

\publ Физматлит

\yr 2004

\totalpages 823

\misctransl

\by Andreev~S.\,G., Babkin~A.\,V., Baum~F.\,A., et al.

\book Physics of Explosion

\ed L.~P.~Orlenko 

\bookvol 1

\publ Fizmatlit

\publaddr Moscow

\yr 2004

\totalpages 823



\RBibitem{reut1:bib9}

\by Кривченко~А.\,Л.

\paper Метод расчёта параметров детонации конденсированных взрывчатых веществ

\jour ФГВ

\yr 1984

\vol 20

\issue 3

\pages 83--86

\transl англ. пер.:~

\jour Combustion, Explosion, and Shock Waves

\vol 20

\issue 3

\pages 324--326

%DOI: 10.1007/BF00782919

\paper Method of computing the detonation parameters of condensed high explosives

\by Krivchenko~A.\,L. 







\RBibitem{reut1:bib3}

\by Казутин~М.\,В., Комаров~В.\,Ф., Попок~Н.\,И.

\paper Степень реализации энергетического потенциала взрывчатого вещества в метательную способность продуктов детонации

\jour Ползуновский вестник

\issue 4--1

\yr 2010

\pages 92--95

\misctransl

\by Kazutin~M.\,V., Komarov~V.\,F., Popok~N.\,I.

\paper The degree of implementation of the energy potential of the explosive in the throwing ability of the detonation products

\jour Polzunovskiy Vestnik

\issue 4--1

\yr 2010

\pages 92--95





\RBibitem{reut1:bib4}

\by Смирнов~С.\,П., Колганов~Е.\,В., Кулакевич~Я.\,С.

\paper Связь метательного действия смесевых и индивидуальных ВВ с их составом и строением

\inbook Современные методы проектирования и отработки ракетно-артиллерийского вооружения

\publaddr Саров

\publ ВНИИЭФ

\yr 2000

\pages 410--412

\misctransl

\by Smirnov~S.\,P., Kolganov~E.\,V., Kulakevich~Ya.\,S.

\paper Relationship of throwing action mixtures and individual explosives to their composition and structure

\inbook Modern methods of design and testing of missile and artillery armament

\publaddr Sarov

\publ VNII\'EF

\yr 2000

\pages 410--412



\RBibitem{reut1:bib5}

\by Пепекин~В.\,И., Губин~С.\,А.

\paper Метательная способность органических взрывчатых веществ и~их пределы по мощности и скорости детонации

\jour ФГВ

\yr 2007

\vol 43

\issue 1

\pages 99--111

\transl англ. пер.:~

\paper Propellant performance of organic explosives and their energy output and detonation velocity limits

\by Pepekin~V.\,I., Gubin~S.\,A.

\jour Combustion, Explosion, and Shock Waves

\yr 2007

\vol 43

\issue 1

\pages 84--95



\Bibitem{reut1:bib6}

\by Makhov~M.\,N.

\paper Explosion Heat and Metal Acceleration Ability of High Explosives

\jour AIP Conf. Proc.

\vol 706

\yr 2004

\pages 863--866

%doi:10.1063/1.1780373 (4 pages)





\Bibitem{reut1:bib7}

\by van Thiel~M., Kusubov~A.\,S., Mitchell~A.\,C.

\preprint Compendium of shock wave data

\preprintinfo  Technical Report UCRL-50108

\publaddr Livermore, CA

\publ Lawrence Radiation Lab.

\yr 1967





\Bibitem{reut1:bib8}

\by Rao~W.\,R.

\paper  Velocity of Sound in Liquids and Chemical Constitution

\jour J. Chem. Phys.

\vol 9

\issue 9

\pages 682--685

\yr 1941

%doi:10.1063/1.1750976 (4 pages)











\RBibitem{reut1:bib10}

\by Лин~Э.\,Э., Пащенко~Э.\,Н., Тихомиров~Б.\,П.

\paper Метание пластин продуктами пересжатой детонации заряда взрывчатого вещества с уменьшающейся плотностью

\jour ПМТФ

\yr 1993

\issue 1

\pages 32--34

\transl англ. пер.:~

\jour J.~Appl. Mech. Tech. Phys.

\vol 34

\issue 1

\pages 28--30

%DOI: 10.1007/BF00851801

\paper Theoretical models of detonation of a flat layer of condensed explosive with diminishing density

\by Lin~\'E.\,\'E., Pashchenko~\'E.\,N., Tikhomirov~B.\,P.





\end{thebibliography}





\noindent

\hskip6.5cm \thedate \par\noindent

\hskip6.5cm\therevdate





\makeengtitlem





 \newpage

%

% \tableofengcontents

%

% \end{document}

